
%Class
	\documentclass[a4paper,12pt]{article}

%Packages
	\usepackage{amsmath}					% Standard maths
	\usepackage{amssymb}					% Standard maths
	\usepackage{amsthm}						% Standard maths
	\usepackage{thmtools, thm-restate}% Theorem styles
	\usepackage{mathtools}
	\usepackage[newzealand]{babel}% Date/time in British format (yes, I know it says NZ)
	\usepackage{datetime}					% Date/time
	\usepackage{multirow}					% For stacked subscripts
	\usepackage{xcolor}						% For colour!
	\usepackage[normalem]{ulem}		% For strikeout
	\usepackage{isomath}
	\usepackage{fancyhdr}					% For the headers
	\usepackage{mfirstuc}					% For the headers
	\usepackage[hmargin=1in,top=0.8in,bottom=1in]{geometry} % Page margins more reasonable
	\usepackage{graphicx}					% For including graphics
	%\usepackage{float}						% For something
	\usepackage{enumitem}					% For numbering tweaks
	\usepackage{pgfplots}
	\usepackage[margin=15pt,font={footnotesize},labelfont=bf,format=hang]{caption}
\usepackage[margin=15pt,font={footnotesize},labelfont=bf,format=hang]{subcaption}
	\usepackage{url}
	\usepackage[hidelinks]{hyperref}	% For hyperlinks
	\usepackage{breakurl}
	\usepackage{cleveref}
	\usepackage[scaled=.92]{helvet}

%MATHEMATICAL SHORTCUTS
%Two and three-part definitions
	\newcommand{\twopartdef}[4]
	{	\left\{
			\begin{array}{ll}
				#1 & \mbox{if } #2 \\
				#3 & \mbox{if } #4
			\end{array}
		\right.}
%	\newcommand{\threepartdef}[6]
%	{	\left\{
%			\begin{array}{ll}
%				#1 & \mbox{if } #2 \\
%				#3 & \mbox{if } #4 \\
%				#5 & \mbox{if } #6
%			\end{array}
%		\right.}
%Blackboard bold	
	\newcommand{\N}{{\mathbb N}} 
	\newcommand{\R}{{\mathbb R}} 
	\newcommand{\C}{{\mathbb C}} 
	%Curly letters
%	\newcommand{\E}{{\mathcal E}}	
%	\newcommand{\F}{{\mathcal F}} 
%	\newcommand{\Null}{{\mathcal N}} 
%	\newcommand{\R}{\mathcal{R}} 
%	\newcommand{\Z}{\mathcal{Z}} 
%Uprights
	\newcommand{\D}{{\mathrm D}}
	\renewcommand{\d}{{\mathrm d}}
%Easy Greeks
%	\def\a{\alpha}
%	\def\g{\gamma}
%	\newcommand{\e}{{\varepsilon}}
%	\newcommand{\la}{{\lambda}}  
%	\newcommand{\om}{{\Omega}}
	\renewcommand{\epsilon}{\varepsilon}
	
%Vector operators
	\def \p {\partial}
	\newcommand{\del}{{\nabla}}
	\newcommand{\grad}{{\nabla}}
    \newcommand{\vdel}{\boldsymbol{\nabla}}
    \newcommand{\vgrad}{\boldsymbol{\nabla}}
    \newcommand{\vnabla}{\boldsymbol{\nabla}}
	\newcommand{\cross}{{\times}}
	\newcommand{\Lap}{{\nabla^2}} 
%Arrows
%	\newcommand{\goesto}[1]{\xrightarrow[#1]{}}
%hats
%	\newcommand{\nhat}{\widehat{\v{n}}}
%	\newcommand{\shat}{\widehat{\v{s}}}
	\renewcommand{\hat}{\widehat}
%Note to self
%	\newcommand{\note}[1]{[\textbf{\textsl{\MakeUppercase{#1}}}]}
	
%Stress tensor
%	\newcommand{\T}{\vg{\sigma}}

% Vectors, quaternions etc.
\renewcommand{\v}[1]{\mathbfit{#1}}      % Generic vector
\newcommand{\vc}[1]{\bm{\mathcal{#1}}}      % vector mathcal
\newcommand{\uv}[1]{\widehat{\mathbfit{#1}}}      % Generic unit vector
\newcommand{\q}[1]{\mathsfbfit{#1}}        % Generic quaternion
\renewcommand{\t}[1]{\mathsfbfit{#1}}	 % Generic tensor
\newcommand{\m}[1]{\mathsfbfit{#1}}	 % Generic matrix
\newcommand{\vu}[1]{\mathbf{#1}}         % Upright vector (for zero vector)
\newcommand{\qu}[1]{\mathbf{#1}}       % Upright quaternion (for zero quaternion)
\newcommand{\tu}[1]{\bm{\mathsf{#1}}}	 % Upright tensor (for zero tensor)

%Real and imaginary
\renewcommand{\Re}{\operatorname{Re}}
\renewcommand{\Im}{\operatorname{Im}}
	
%Non-dim numbers
%	\renewcommand{\Re}{\operatorname{Re}}
%	\newcommand{\Bi}{\operatorname{Bi}}
%	\newcommand{\Fo}{\operatorname{Fo}}
	
%Misc
\DeclareMathSymbol{\Delta}{\mathalpha}{operators}{1} % Keeps Delta upright
\DeclareMathSymbol{\Gamma}{\mathalpha}{operators}{0} % Keeps Gamma upright
	\newcommand{\cosec}{\operatorname{cosec}}
	\newcommand{\cosech}{\operatorname{cosech}}
	\newcommand{\sech}{\operatorname{sech}}
	\newcommand{\tr}{\operatorname{tr}}
	\newcommand{\erf}{\operatorname{erf}}
	\newcommand{\order}{\mathcal{O}}
%	\newcommand{\V}{\mathcal{V}} % 	CurlyV = (Bi V) / (1 + Bi)

% Differentiating 
\newcommand{\fd}[2]{\mathchoice{\frac{\d #1}{\d #2}}{\d #1/\d #2}{\d #1/\d #2}{\d #1/\d #2}}
\newcommand{\fdd}[2]{\mathchoice{\frac{\d^2 #1}{\d #2^2}}{\d^2 #1/\d #2^2}{\d^2 #1/\d #2^2}{\d^2 #1/\d #2^2}}
\newcommand{\fddd}[2]{\mathchoice{\frac{\d^3 #1}{\d #2^3}}{\d^3 #1/\d #2^3}{\d^3 #1/\d #2^3}{\d^3 #1/\d #2^3}}
\newcommand{\fdddd}[2]{\mathchoice{\frac{\d^4 #1}{\d #2^4}}{\d^4 #1/\d #2^4}{\d^4 #1/\d #2^4}{\d^4 #1/\d #2^4}}
\newcommand{\fdn}[3]{\mathchoice{\frac{\d^{#1} #2}{\d #3^{#1}}}{\d^{#1} #2/\d #3^{#1}}{\d^{#1} #2/\d #3^{#1}}{\d^{#1} #2/\d #3^{#1}}}
\newcommand{\pd}[2]{\mathchoice{\frac{\p #1}{\p #2}}{\p #1/\p #2}{\p #1/\p #2}{\p #1/\p #2}}
\newcommand{\pdd}[2]{\mathchoice{\frac{\p^2 #1}{\p #2^2}}{\p^2 #1/\p #2^2}{\p^2 #1/\p #2^2}{\p^2 #1/\p #2^2}}
\newcommand{\pddmixed}[3]{\frac{\p^2 #1}{\p #2 \p #3}}
\newcommand{\pddd}[2]{\frac{\p^3 #1}{\p #2^3}}
\newcommand{\pdn}[3]{\mathchoice{\frac{\p^{#1} #2}{\p #3^{#1}}}{\p^{#1} #2/\p #3^{#1}}{\p^{#1} #2/\p #3^{#1}}{\p^{#1} #2/\p #3^{#1}}}
	
	% Misc
\newcommand{\e}{\mathrm{e}}
\renewcommand{\i}{\mathrm{i}}
\newcommand{\dx}{\,\mathrm{d}x}
\newcommand{\dy}{\,\mathrm{d}y}
\renewcommand{\epsilon}{\varepsilon}
\renewcommand{\Re}{\operatorname{Re}}
\newcommand{\Four}{\mathcal{F}}

\newcommand{\mat}[4]{\begin{pmatrix}#1 & #2 \\ #3 & #4 \end{pmatrix}}

	%Column vectors
	\makeatletter
\newcommand\rcvector[2][\\]{\ensuremath{%
  \global\def\rc@delim{#1}%
    \negthinspace\begin{pmatrix}
      \rc@vector #2,\relax\noexpand\@eolst%
    \end{pmatrix}}}
\def\rc@vector #1,#2\@eolst{%
  \ifx\relax#2\relax
    #1
  \else
    #1\rc@delim
    \rc@vector #2\@eolst%
  \fi}
\makeatother

\newcommand{\colvec}{\rcvector}
\newcommand{\rowvec}{\rcvector}
\renewcommand{\binom}{\rcvector}
	
	%Colours
%	\newcommand{\orange}[1]{{\color{orange} #1 }}
%	\newcommand{\blue}[1]{{\color{blue} #1 }}

%Theorem styles
%	\declaretheoremstyle[headpunct={\:\:\:}, shaded={bgcolor={rgb}{0.9,0.9,0.9}, margin=5pt}]{problem}
%	\declaretheoremstyle[headpunct={\:\:\:}, shaded={rulecolor=black, rulewidth=0.4pt, bgcolor={rgb}{1,1,1}, margin=5pt}]{definition}
%	\declaretheoremstyle[headpunct={\:\:\:}, spaceabove=15pt, notefont=\bf, notebraces={(}{)}]{theorem}
%	\declaretheoremstyle[headpunct={\:\:\:}, numbered=no, spaceabove=15pt, qed={\checkmark}]{solution}
%	
%	\declaretheorem[style=definition,numberwithin=section]{definition}
%	\declaretheorem[numberlike=definition, style=theorem]{theorem}
%	\declaretheorem[numberlike=definition, style=theorem]{lemma}
%	\declaretheorem[numberlike=definition,style=problem]{problem}
%	\declaretheorem[numberlike=definition, style=theorem]{proposition}
%	\declaretheorem[numberlike=definition, style=theorem]{corollary}
%	\declaretheorem[numberlike=definition, style=theorem]{remark}
%	\declaretheorem[numberlike=definition, style=theorem]{example}
%	\declaretheorem[numberlike=definition, style=theorem]{notation}
%	\declaretheorem[style=solution]{solution}

%Depth of display breaks to allow	
	\allowdisplaybreaks[1]

%How deep do we want the Table of Contents to go?	
%	\setcounter{tocdepth}{1}

%Sections
\makeatletter
\renewcommand{\section}{\@startsection
{section}%                   % the name
{1}%                         % the level
{0mm}%                       % the indent
{-\baselineskip}%            % the before skip
{0.5\baselineskip}%          % the after skip
{\normalfont\bfseries}} % the style
\makeatother

% Wavy underline
\makeatletter
\font\sixly=lasyb10 scaled 1200
\def\tinners{\bgroup \markoverwith{\lower8.5\p@\hbox{\sixly \char58}}\ULon}
\makeatother
\newcommand{\answer}[1]{\tinners{\; #1 \;}}

%Setting the footnotes to use *,**,+ etc rather than 1,2,3	
\renewcommand{\thefootnote}{\fnsymbol{footnote}}

%Italic quote environment
%\newenvironment{quoteitalic}{\begin{quote}\itshape}{\end{quote}}

%Heading
\newcommand{\heading}[1]{
\setlength{\parindent}{0pt}
\textbf{MATH3171 (Mathematical Biology III)}

\textbf{YEAR 2025--26}%, MICHAELMAS TERM}
\setlength{\parskip}{3ex}

\textbf{\MakeUppercase{#1}}
%
\setlength{\parskip}{2ex}

}

%========================================================================================================================
%DOCUMENT BEGINS HERE!

\usepackage{multicol}
\setlength{\columnsep}{3pc}
\setlength\columnseprule{.4pt}
\begin{document} 
%
\heading{Past paper guide}

This guide suggests how well previous years' exam questions fit to this year's course. Questions fall into the following categories:
\begin{itemize}
	\item A \textbf{good fit} means an exam question could look like this.
	\item A \textbf{reasonable fit} means the mathematical techniques are relevant but the style is different. A such, it is worth doing if you have exhausted all the `good fit' questions.
	\item If a question is not listed it is \textbf{not relevant} to this year's course. If specific parts aren't mentioned, then all of the question is relevant.
\end{itemize}

Terminology changes a bit throughout the papers but know that
\begin{itemize}
	\item Critical point = steady state = fixed point = equilibrium.
	\item Constant non-zero steady state = spatially homogeneous equilibrium = homogeneous equilibrium.
	\item Physically meaningful equilibrium = permissible equilibrium = feasible equilibrium.
	\item Classify the equilibrium\dots = what type of equilibrium\dots = decide whether it is a stable node etc.
\end{itemize}

\vspace{1cm}

\begin{multicols}{2}

\section*{2025 paper}
\begin{itemize}
\item[Q1] Good fit.
\item[Q2] Reasonable fit.
\item[Q3] Good fit.
\item[Q4] Good fit.
\item[Q5] Good fit.
\item[Q6] Reasonable fit.
\item[Q7] Good fit.
\item[Q8] Good fit.
\end{itemize}

\section*{2024 paper}
\begin{itemize}
\item[Q1] Good fit.
\item[Q2] Good fit.
\item[Q3] Good fit.
\item[Q4] Good fit.
\item[Q5] Reasonable fit.
\item[Q6] (a)--(c) Good fit. (d) is subtly different to the cases we've seen in class because there is no growth term, which means all $u$ are homogeneous equilibria and the boundary conditions are different.
\item[Q7] Good fit.
\item[Q8] Good fit.
\end{itemize}

\section*{2023 paper}
\begin{itemize}
\item[Q1] Good fit.
\item[Q2] Good fit.
\item[Q3] Good fit.
\item[Q4] Good fit.
\item[Q5] Good fit.
\item[Q6] Good fit.
\item[Q7] Good fit.
\item[Q8] Good fit.
\end{itemize}

\section*{2022 paper}
\begin{itemize}
\item[Q1] Good fit.
\item[Q2] Good fit.
\item[Q3] Good fit.
\item[Q4] Good fit.
\item[Q5] Good fit.
\item[Q6] Good fit.
\item[Q7] Good fit.
\item[Q8] Good fit.
\end{itemize}

\section*{2021 paper}
Take-home paper, so questions are longer than you would expect this year.
\begin{itemize}
\item[Q1] Good fit.
\item[Q2] Good fit.
\item[Q3] Good fit.
\item[Q4] Good fit.
\end{itemize}

\section*{2020 paper}
Take-home paper, so questions are longer than you would expect this year.
\begin{itemize}
\item[Q1] Good fit.
\item[Q3] Good fit.
\item[Q4] Good fit.
\item[Q5] Good fit.
\end{itemize}

\section*{2019 paper}
\begin{itemize}
\item[Q1] Good fit.
\item[Q2] Good fit.
\item[Q3] Good fit.
\item[Q4] Good fit.
\item[Q5] Reasonable fit.
\item[Q6] Good fit.
\item[Q7] Good fit.
\item[Q8] Good fit. 
\item[Q9] Good fit.
\end{itemize}

\section*{2018 paper}
\begin{itemize}
\item[Q1] Good fit.
\item[Q2] Good fit.
\item[Q3] Good fit.
\item[Q4] Good fit.
\item[Q5] Good fit.
\item[Q6] Good fit.
\item[Q7] Good fit.
\item[Q8] Good fit. 
\item[Q10] Reasonable fit.
\end{itemize}

\section*{2017 paper}
\begin{itemize}
\item[Q1] Good fit.
\item[Q2] Good fit.
\item[Q3] Reasonable fit.
\item[Q4] Good fit.
\item[Q5] Good fit.
\item[Q6] Good fit.
\item[Q7] Good fit.
\item[Q8] Good fit. 
\end{itemize}

\section*{2016 paper}
\begin{itemize}
\item[Q1] Good fit.
\item[Q2] Good fit.
\item[Q3] Good fit.
\item[Q5] Good fit.
\item[Q6] Good fit.
\item[Q7] Good fit.
\item[Q8] Good fit. In (e), conditions for aggregation mean satisfying the Turing conditions.
\item[Q9] Good fit.
\item[Q10] Reasonable fit. In (d), the linear expansion is $u=u_0 + \delta u$.
\end{itemize}

\section*{2015 paper}
\begin{itemize}
\item[Q1] Reasonable fit. 
\item[Q3] (a)--(b) Good fit. 
\item[Q4] Good fit. Linear (instability) analysis = linear stability analysis.
\item[Q6] Good fit.
\item[Q7] (a)--(e) Reasonable fit.
\item[Q8] Good fit.
\item[Q9] Reasonable fit.
\item[Q10] Good fit, but they would have seen this model already in class.
\end{itemize}

\section*{2014 paper}
\begin{itemize}
\item[Q1] Good fit. The definition of $J$ is such that $\pd{\rho}{t} = -\pd{J}{x}$. You would be given $J$ in an exam.
\item[Q3] Good fit.
\item[Q4] Good fit.
\item[Q5] Good fit. 
\item[Q6] Good fit.
\item[Q7] Good fit. In (c), the linear expansion is $u = u^* + v$ where $v$ is small (so these are the $\mathcal{O}(v)$ equations).
\item[Q8] Apart from (d), reasonable fit.
\item[Q10] Good fit. In (d), conditions for aggregation mean satisfying the Turing conditions.
\end{itemize}

\section*{2013 paper}
\begin{itemize}
\item[Q1] (a), (b) Good fit. (c), (d) Reasonable fit.
\item[Q2] Good fit.
\item[Q4] Good fit.
\item[Q5] Good fit. 
\item[Q6] Reasonable fit.
\item[Q7] Good fit. The definition of $J$ is such that $\pd{c}{t} = -\pd{J}{x}$ and the mean transit time is $N/J$.
\item[Q8] Reasonable fit.
\item[Q9] Good fit.
\item[Q10] Reasonable fit. Note that $\Delta = \nabla^2$, the Laplacian. In (d) the perturbation is in the form $\rho = \rho_0 + \delta \rho$ where $\delta \rho$ is small. In (f) aggregation implies the Turing instability. Aggregation comes from requiring the instability by $\sigma>0$.
\end{itemize}

\section*{2012 paper}
\begin{itemize}
\item[Q2] (a)--(c) Good fit.
\item[Q3] (a), (b) Reasonable fit.
\item[Q4] Good fit.
\item[Q5] Good fit. 
\item[Q8] Good fit, but they would have seen this model already in class.
\item[Q9] Reasonable fit.
\item[Q10] Reasonable fit.
\end{itemize}

\section*{2011 paper}
\begin{itemize}
\item[Q1] Good fit
\item[Q2] Reasonable fit. Not sure if this is how Shakespeare meant it, mind you.
\item[Q3] Good fit. 
\item[Q4] Reasonable fit.
\item[Q5] Good fit.
\item[Q6] Good fit. For (d) two possibilities should be enough. Chemotactic wave = travelling wave.
\item[Q7] Good fit. Use Fourier transforms to solve this directly (I think here they are looking for you to show it is true by differentiation.)
\item[Q8] Reasonable fit.
\item[Q9] Reasonable fit.
\item[Q10] Good fit.
\end{itemize}

\end{multicols}
%
\end{document}
